\subsection{Epoch-budget validation}
\label{sec:phase4}

Every result in Sections~\ref{sec:kan-vs-mlp}--\ref{sec:track-e} was measured
at a fixed $10^4$-epoch budget, a tenth of the base paper's, for the cost
reasons of Section~\ref{sec:budget}. Here we ask whether that budget is long
enough for the \emph{rankings} to be trusted, and whether any gap to the
paper comes from the budget or from an implementation difference.

\begin{expectbox}
\expectrow{Question}{Do the conclusions drawn at $10^4$ epochs survive five times the training budget?}
\expectrow{Expected}{Absolute errors fall; orderings between configurations hold, since every $10^4$-epoch model had already reached a training loss near $10^{-4}$.}
\expectrow{Observed}{Three of four conclusions hold. The fourth, the one the paper headlines, reverses: at $5\times10^4$ epochs the parameter-matched MLP extrapolates $6.8\times$ \emph{better} than KAN-ODE.}
\end{expectbox}

\subsubsection{Design}

We chose five configurations to test the conclusions most exposed to the
budget: Euler (the one solver short of $10^{-4}$ at $10^4$ epochs),
Tsit5/RBF (the reference), B-spline (the near-tie with RBF), and both MLP
baselines. Each was retrained from scratch at $5\times10^4$ epochs with the
recipe of Table~\ref{tab:recipe} unchanged, on cloud CPUs. The B-spline run
was cut to $2.5\times10^4$ epochs because its short calibration run
predicted about 35 hours for the full budget. The real run turned out
$1.63\times$ faster than that calibration ($1.544$ against $2.52$~s/epoch),
a reminder that a five-epoch sample is dominated by start-up cost. The batch
used $23.2$ serial CPU-hours against a planned $17.0$, and $10.7$ hours of
wall-clock with the five jobs in parallel. Before analysis we checked every
run's stored metrics against the epoch count of its training log, and
reloaded and re-scored every checkpoint; the rollouts reproduce the stored
extrapolation MSE to within $0.6\%$.

We added a sixth configuration later, as a separate follow-up. The
implementation audit (Section~\ref{sec:audit}) turned up the base paper's own
learning rate and initialization scale for its MLP baseline, so we reran
both budgets with the identical $[2,50,2]$ tanh architecture under those
settings instead of our default. We report it alongside the original five
throughout.

\subsubsection{Results}

\begin{table*}[t]
\centering
\footnotesize
\renewcommand{\arraystretch}{1.1}
\caption{The five originally planned configurations at $10^4$ epochs
(Sections~\ref{sec:kan-vs-mlp}--\ref{sec:basis-results}) and at the extended budget, plus a sixth added
after auditing the paper's own optimizer settings for its MLP baseline
(Section~\ref{sec:audit}). ``Clip'' is the fraction of optimizer steps on
which gradient clipping at norm $1.0$ engaged; $\|\nabla\|_{\text{med}}$ is
the median pre-clip gradient norm over the extended run. The sixth
configuration ran with no clipping by design (see the paragraph following
this table), so those two columns are blank throughout.}
\label{tab:phase4}
\begin{tabularx}{\textwidth}{@{}Y C{1.1cm} C{2.0cm} C{2.0cm} C{1.8cm} C{1.0cm} C{1.1cm} C{1.6cm}@{}}
\toprule
\hdrrow \thh{Configuration} & \thh{Epochs} & \thh{Training MSE} & \thh{Extrap.\ MSE} & \thh{Extrap.\ $R^2$} & \thh{$\hat L$} & \thh{Clip} & \thh{$\|\nabla\|_{\text{med}}$}\\
\midrule
Euler / RBF & 10k & $1.38\times10^{-4}$ & $2.51\times10^{-3}$ & $0.99914$ & $8.55$ & & \\
 & \textbf{50k} & $8.85\times10^{-6}$ & $1.93\times10^{-4}$ & $0.99993$ & $7.72$ & $4.7\%$ & $0.147$\\
\addlinespace
Tsit5 / RBF & 10k & $8.83\times10^{-5}$ & $8.92\times10^{-5}$ & $0.99997$ & $6.91$ & & \\
 & \textbf{50k} & $6.59\times10^{-6}$ & $4.02\times10^{-5}$ & $0.999986$ & $7.75$ & $4.3\%$ & $0.169$\\
\addlinespace
Tsit5 / B-spline & 10k & $7.20\times10^{-5}$ & $8.38\times10^{-5}$ & $0.99997$ & $8.44$ & & \\
 & \textbf{25k} & $9.65\times10^{-6}$ & $4.26\times10^{-5}$ & $0.999985$ & $7.05$ & $2.2\%$ & $0.033$\\
\addlinespace
MLP-ODE, SiLU & 10k & $9.50\times10^{-5}$ & $4.77\times10^{-4}$ & $0.99984$ & $5.07$ & & \\
 & \textbf{50k} & $\mathbf{1.82\times10^{-6}}$ & $\mathbf{5.90\times10^{-6}}$ & $\mathbf{0.999998}$ & $5.65$ & $11.1\%$ & $0.162$\\
\addlinespace
MLP-ODE, tanh (paper, default lr/init) & 10k & $1.08\times10^{0}$ & $1.69\times10^{0}$ & $0.4215$ & $38.2$ & & \\
 & \textbf{50k} & $1.01\times10^{0}$ & $4.32\times10^{0}$ & $-0.4841$ & $104.8$ & $99.3\%$ & $19.8$\\
\addlinespace
MLP-ODE, tanh (paper's own lr/init) & 10k & $1.11\times10^{-3}$ & $9.65\times10^{-3}$ & $0.99669$ & $9.23$ & & \\
 & \textbf{50k} & $5.43\times10^{-5}$ & $4.31\times10^{-4}$ & $0.99985$ & $9.67$ & -- & -- \\
\bottomrule
\end{tabularx}
\end{table*}

\begin{figure*}[t]
\centering
\begin{subfigure}[b]{0.32\textwidth}
  \includegraphics[width=\linewidth]{figures/phase4/tsit5_rbf_50k_loss_curves.png}
  \caption{KAN-ODE, Tsit5 / RBF, 50k}
\end{subfigure}\hfill
\begin{subfigure}[b]{0.32\textwidth}
  \includegraphics[width=\linewidth]{figures/phase4/mlp_silu_50k_loss_curves.png}
  \caption{MLP-ODE, SiLU, 50k}
\end{subfigure}\hfill
\begin{subfigure}[b]{0.32\textwidth}
  \includegraphics[width=\linewidth]{figures/phase4/euler_50k_loss_curves.png}
  \caption{KAN-ODE, Euler / RBF, 50k}
\end{subfigure}\\[4pt]
\begin{subfigure}[b]{0.32\textwidth}
  \includegraphics[width=\linewidth]{figures/phase4/bspline_25k_loss_curves.png}
  \caption{KAN-ODE, Tsit5 / B-spline, 25k}
\end{subfigure}\hfill
\begin{subfigure}[b]{0.32\textwidth}
  \includegraphics[width=\linewidth]{figures/phase4/mlp_paperspec_50k_loss_curves.png}
  \caption{MLP-ODE, tanh (paper's default lr/init), 50k}
\end{subfigure}\hfill
\begin{subfigure}[b]{0.32\textwidth}
  \includegraphics[width=\linewidth]{figures/phase4/mlp_tanh_exact_50k_loss_curves.png}
  \caption{MLP-ODE, tanh (paper's own lr/init), 50k}
\end{subfigure}
\caption{Training and test loss of the six extended runs, as saved by each
run. Training loss is logged every epoch; the test loss (full horizon
$[0,14]$) every ten epochs, and it spikes by up to two orders of magnitude
throughout, which is why the comparison in the text uses a 500-epoch rolling
median rather than raw values. Panel (f) shows the identical $[2,50,2]$ tanh
architecture of panel (e), retrained with only the learning rate and
initialization scale changed to the values in the paper's own published
code.}
\label{fig:phase4-curves}
\end{figure*}

Every working configuration improves its training MSE by $7.5$--$52\times$
(Table~\ref{tab:phase4}); what matters is which orderings survive.
Table~\ref{tab:phase4-verdicts} gives the verdict for each conclusion the
earlier sections draw.

\begin{table*}[t]
\centering
\footnotesize
\renewcommand{\arraystretch}{1.15}
\caption{Which $10^4$-epoch conclusions survive the extended budget.}
\label{tab:phase4-verdicts}
\begin{tabularx}{\textwidth}{@{}L{4.6cm} L{3.2cm} L{3.0cm} Y@{}}
\toprule
\hdrrow \thh{Conclusion at $10^4$ epochs} & \thh{At $10^4$} & \thh{At $5\times10^4$} & \thh{Verdict}\\
\midrule
RBF and B-spline are statistically tied & $6.1\%$ apart & $5.6\%$ apart (25k) & \textbf{Holds}. The basis choice can be made on cost.\\
Euler is the worst solver & $28.1\times$ Tsit5 & $4.8\times$ Tsit5 & \textbf{Holds}, with a much smaller margin: part of the gap was under-training.\\
The paper-spec MLP fails because of its architecture, not its budget & train $1.08$, $R^2=0.42$ & train $1.01$, $R^2=-0.48$ & \textbf{Budget part holds; architecture part reversed.} More epochs alone do not help, but the identical architecture trains under the paper's own learning rate and initialization (below): $R^2=0.9998$ at this budget.\\
KAN-ODE extrapolates better than a parameter-matched MLP & KAN $5.35\times$ better & MLP $6.82\times$ better & \textbf{Reverses.} The advantage is a budget effect.\\
\bottomrule
\end{tabularx}
\end{table*}

\paragraph{The paper-spec MLP is stuck, not slow.} With five times the
budget its training loss moves by $7\%$ ($1.083\to1.011$), its extrapolation
error rises $2.6\times$, and $R^2$ turns negative. The optimizer spends the
run in a pathological regime: clipping engages on $99.3\%$ of steps against
$2$--$11\%$ for every other configuration, the median pre-clip gradient norm
is $19.8$ against $0.03$--$0.17$, and the Lipschitz estimate nearly triples
to $104.8$, an order of magnitude above every other model. This settles half
of the question left open in Section~\ref{sec:kan-vs-mlp}: whatever is wrong
here is a property of this training run, not of the epoch count, since five
times the budget only entrenches it. It is not, however, a property of the
architecture itself. The next paragraph shows the identical network
training cleanly once its optimizer settings are corrected.

\paragraph{Resolving the audit's open optimizer row: the paper's own learning
rate and initialization.} The base paper's own published code (not
distributed with our project, but readable in its public repository) sets
two things differently for its MLP baseline than our default: a learning
rate of $10^{-2}$ (five times our $2\times10^{-3}$) and Glorot initial
weights divided by $10^5$, near zero rather than full-scale. Holding
architecture, activation, dataset, solver and seed fixed and changing only
those two settings, the identical $[2,50,2]$ tanh network trains cleanly at
both budgets (Table~\ref{tab:phase4}, row six; panel f of
Figure~\ref{fig:phase4-curves}). At $10^4$ epochs, training MSE falls from
$1.08$ to $1.11\times10^{-3}$ and extrapolation MSE from $1.69$ to
$9.65\times10^{-3}$ ($R^2:0.42\to0.997$). At $5\times10^4$ epochs training
MSE reaches $5.43\times10^{-5}$ and extrapolation MSE $4.31\times10^{-4}$
($R^2=0.99985$), against $1.01$ and $4.32$ under the default settings at the
same budget. Neither run needed gradient clipping or hit a single
non-finite gradient step. The conclusion Table~\ref{tab:mlp-breakdown} drew
does not survive this: changing only the learning rate and initialization
scale, with the exact literal architecture and activation that table showed
failing, is enough on its own. What that factorial actually measured was
sensitivity to activation \emph{at our optimizer settings}, not an
architecture-level incompatibility with the system.

Against the other two architectures at $5\times10^4$ epochs, the corrected
tanh MLP is still not competitive. Its extrapolation MSE
($4.31\times10^{-4}$) is $10.7\times$ KAN-ODE's and $73\times$ the SiLU MLP's,
and its learned equilibrium sits $1.72$ from the truth, the furthest of the
three (Table~\ref{tab:equilibria}). Its one advantage is conservation: a
first-integral drift of $0.06\%$ is the smallest of the three (KAN $0.08\%$,
SiLU-MLP $0.21\%$), though its on-orbit field error ($2.7\%$) is the
largest. It also falls short of the paper's own reported target. The paper
states that its MLP reaches a training loss of $3.0\times10^{-5}$ within
$10^5$ epochs; this run's training loss crosses $10^{-4}$ at epoch
$39{,}695$ but had not reached $3.0\times10^{-5}$ by epoch $50{,}000$. That
trend is consistent with, but does not confirm, the paper's own budget for
this baseline.

\paragraph{The KAN advantage over the MLP reverses.} At $10^4$ epochs, from
training losses $7\%$ apart, the KAN-ODE extrapolates $5.35\times$ better
than the SiLU MLP, the result that reproduces the paper's accuracy claim. At
$5\times10^4$ epochs the MLP's extrapolation MSE is $5.90\times10^{-6}$ and
the KAN's $4.02\times10^{-5}$: the MLP is $6.82\times$ better, far outside
the single-seed noise that separates RBF from B-spline. The checkpoint
analyses of Section~\ref{sec:numerics} agree on every measure. At 50k the MLP
has the smaller on-orbit field error ($0.6\%$ against $1.1\%$), the smaller
first-integral drift ($0.4\%$ against $0.9\%$), an equilibrium closer to the
truth ($0.68$ against $1.65$ away), and a lower far-horizon error
(maximum $2.4\times10^{-2}$ against $6.5\times10^{-2}$ on $(14,28]$). The
crossover is late. Smoothing the logged full-horizon MSE of
Figure~\ref{fig:phase4-curves} with a 500-epoch rolling geometric median,
the two models are level at $2.5\times10^4$
epochs ($1.7\times10^{-4}$ against $2.0\times10^{-4}$), and the MLP stays
ahead for good only from about epoch $48{,}500$. The MLP is also the faster
learner at every threshold: in the extended runs it reaches $10^{-4}$
training MSE at epoch 5225 against the KAN's 5695, and $10^{-5}$ at 13{,}815
against 31{,}412.

The honest reading is that KAN-ODE has a strong \emph{early} inductive bias
for this smooth periodic problem. Within a fixed, moderate budget it produces
the better extrapolator, but it is not the better function class once both
models are trained to convergence. Neither model learns the conservative
structure of the system (Figure~\ref{fig:limit-cycle}); the MLP simply fits
the training orbit's phase more precisely in the end.
